\chapter{Characterization of the phase plate}
The final goal of this project was to have a complete characterization of the phase plate that we received from the company. We knew its expected behaviour, namely the creation of a top hat potential with the characteristics described in \Cref{sec:slm_phaseplate}. However, it is important to test any new components, to see how they behave in reality.

\section{Building the experimental setup}
\subsection{Imaging system}
In order to characterize the phase plate, the first step was to set up an imaging system. This was necessary to acquire images of the intensity distribution generated by our spatial light modulating system. Since the expected dimensions of the intensity distribution were of the order of \SI{10}{\micro\meter}, a magnification system was set up. The magnification system is composed by a microscope objective, which magnifies the beam, and by a CCD camera. The distance between the two is kept constant by fixing everything in a cage. The objective can be aligned to the beam moving it in the horizontal ($y$) and vertical ($z$) direction by turning two micrometers. A sketch of the setup is shown in \cref{fig:imaging}. The whole system (objective+camera) was mounted on a linear translation stage, with which it could be moved along the $x$ direction to find the focus.

\begin{figure}
    \centering
    \includegraphics[width=0.7\textwidth]{chapters/chapter_3/figures/imaging}
    \caption{Imaging system. A microscope objective magnifies the image and a CCD camera acquires it. The objective can be moved in the $y$ and $z$ direction with two micrometers.}
    \label{fig:imaging}
\end{figure}

In order to use the imaging system to take measurements, it had to be calibrated. For the calibration, a laser beam was sent to a target like the one shown in \cref{fig:target} and imaged. From the image shown in \cref{fig:magnified}, it was possible to calculate the magnifying power of the imaging system. The lines were measured to be separated by \SI{413(2)}{\micro\meter}. Knowing that on the target the separation was $1/55$~mm, we find a magnifying power of \SI{22.7(1)}{}.

\begin{figure}
    \begin{subfigure}[b]{0.35 \textwidth}
        \includegraphics[width=\textwidth]{chapters/chapter_3/figures/target.jpg}
        \\
        \caption{Target}
        \label{fig:target}
    \end{subfigure}
    \hfill
    \begin{subfigure}[b]{0.6\textwidth}
        \includegraphics[width=\textwidth]{chapters/chapter_3/figures/magnified}
        \caption{Magnified image}
        \label{fig:magnified}
    \end{subfigure}
    \caption{Calibration of the imaging system. On the left, an example of a target similar to the one used in the experiment. On the right, the magnified image of a portion of the target. The number 55 means that the target has 55 lines / \si{mm}.}
\end{figure}

\subsection{Test with the $0-\pi$ phase plate}
Before testing the top hat phase plate, the setup was tried with the $0-\pi$ phase plate currently used. In this previous setup, the laser beam coming from a fiber is collimated with a single lens. The beam is then magnified in the $z$ direction by a telescope made of two cylindrical lenses ($f=-40/150$~mm) placed \SI{110}{mm} apart. The magnification of the beam helps to achieve higher trapping frequencies. The beam is then sent to the phase plate and an $f=\SI{100}{mm}$ lens. The two cylindrical lenses and the phase plate are mounted on a rotation mount, to allow the alignment of their axis. The result is imaged by the imaging system at the focal plane of the last lens. The focus is adjusted by moving the imaging system along the $x$ direction with a micrometer. A sketch of this setup is shown in \cref{fig:0pi_setup}.

\begin{figure}
    \includegraphics[width=0.9\textwidth]{chapters/chapter_3/figures/0pi_setup.pdf}
    \caption{Setup for the test of the $0-\pi$ phase plate. The beam coming from the fiber is collimated by a single lens, then expanded by two cylindrical lenses and sent to the phase plate. After the phase plate, an $f=\SI{100}{mm}$ lens is placed and the result is imagined at the focal plane of the last lens. The two cylindrical lenses and the phase plate are mounted on a rotation mount.}
    \label{fig:0pi_setup}
\end{figure}

In order to get a good result, it is important that the two cylindrical lenses and the phase plate axes are aligned. To align it, the beam was observed at short and long distances. Acting on the first lens, the beam was rotated in a way such that its axis was aligned vertically. The same alignment was then performed acting on the second cylindrical lens and looking at the beam in the far field. The procedure was repeated, \enquote{walking} the beam, until the two lenses were aligned. The collimation of the beam was ensured by using a shear plate. Finally, the phase plate was aligned rotating its axis in such a way that it was parallel to the beam axis and moving it up or down such that the intensity of the two peaks was the same.

The final result is shown in \cref{fig:0pi_image}. The result resembles the expected $\text{TEM}_{01}$ mode shown in \cref{fig:tem10}, but it is clear that it could be optimized and made more symmetric. In particular, the alignment of the whole setup could have been improved.  However, since the characterization of the $0-\pi$ phase plate was not the goal of the project, the setup was not optimized. Nonetheless, this test was useful to gain familiarity on how to work with phase plates, before studying the one we were interested in.

\begin{figure}
    \centering
    \includegraphics[width=0.9\textwidth]{chapters/chapter_3/figures/0pi_image.pdf}
    \caption{Intensity distribution after the $0-\pi$ phase plate. The size shown on the axis is the real size of the beam. The plots on the right and on top of the image show the integrated sum of the intensity along the respective axes.}
    \label{fig:0pi_image}
\end{figure}

\subsection{Shaping the input beam}

